\section{Introducción}
\label{iv:introduccion}

La relación entre una estructura matemática y su realización física
requiere conocer qué operaciones producen el objeto, qué información
conservan sus representaciones y qué ecuaciones se transportan entre
ellas. Una dimensión, un valor espectral o un grupo de automorfismos
constituyen resultados distintos de esa investigación. Cuando proceden
de un mismo estado, su relación se expresa mediante los mapas que los
producen. Este artículo estudia esa relación en la prolongación de HMT
hacia la incidencia excepcional, las realizaciones hilbertianas y los
sistemas de compactificación.

La tesis que organiza la serie sitúa la estructura discreta del continuo
antes de su evaluación real. APP proporciona las dos evaluaciones
aritméticas sobre un soporte común; TRIT conserva los regímenes
orientados; TPK selecciona, transporta y prolonga los estados con sus
datos de frontera y memoria. La reducción de cilindros compatibles
determina coordenadas numéricas. La proyección excepcional conserva la
incidencia que se realiza mediante códigos, retículos y álgebras
graduadas. El valor y la estructura que lo genera permanecen enlazados
por esta procedencia, aun cuando cada representación retenga información
diferente.

La función del registro \(K\) permite seguir ese enlace de forma
especialmente concreta. Primero se construyen sus doce componentes
desde el registro firmado y las operaciones del TPK. Su lectura
racional conserva el orden de los bloques con base y origen de fase
fijados. La incidencia de esos componentes interviene después en la
bandera que conduce del diseño de Witt al retículo de Leech.
En otra representación declarada, el mismo registro produce la
dirección \(u_K\) que distingue una recta dentro del espacio
\(V_{11}\). La reducción a \(V_{10}\) tiene entonces un núcleo
especificado, una norma exacta y una ley de covariancia. Este recorrido
explica por qué es indispensable conservar la construcción de \(K\)
antes de utilizar su valor o su dirección.

El corredor excepcional se expone hasta \(V^\natural\), con la
construcción de Frenkel--Lepowsky--Meurman y los resultados de Moonshine
en sus dominios propios~\cite{flm,borcherds}. La elección del retículo,
la composición de sus sectores y la acción de sus automorfismos ocupan
lugares distintos dentro de la prueba. Las realizaciones de órdenes
\(2\) y \(3\) se incorporan con sus referencias y transformaciones.
La función modular se relaciona posteriormente con la inversión del
radio mediante una identidad explícita en el eje imaginario del
semiplano superior. Esa identidad permite componer una relación precisa
entre parámetros sin sustituir los operadores de cada representación.

La realización hilbertiana aporta otro paso necesario. En cada
profundidad se conserva la etiqueta de fase de una base ortonormal.
El desplazamiento y el carácter de fase generan toda el álgebra
matricial de esa fibra. Al refinar, la inclusión mediante la identidad
de la fibra nueva conserva la unidad y la traza normalizada; fijar una
sola continuación produce una inclusión en una esquina. Las pruebas
del cierre inductivo y de la realización tracial hacen visible la
consecuencia de esa elección. La conservación de la unidad deja así
de ser sólo una formulación conceptual: interviene en la acción de
las inclusiones y en el tipo del álgebra resultante.

La dualidad compacta se desarrolla a continuación sobre un dominio
estable. El intercambio de las coordenadas enteras exige transportar
también el radio y, cuando cambia la sección residual, la relación de
equivalencia que conserva el entero. El artículo demuestra la
involución, la invariancia de la forma cuadrática, la autoadjunción
del Hamiltoniano y la igualdad de sus dominios bajo conjugación
unitaria. La evaluación sobre clases incluye el cociente euclídeo;
su mínimo puede calcularse mediante dos candidatos enteros. Así se
conserva el contenido aritmético del cambio de representante durante
la construcción de la equivalencia espectral.

Las pantallas dimensionales reúnen estos resultados con la incidencia.
La perforación de un código, la eliminación ortogonal de una dirección
y el cociente isotrópico de un retículo son operaciones diferentes.
Cada una determina un rango y conserva un núcleo propio. Su
coordinación utiliza esas presentaciones completas. La realización
undecadimensional recibe después las pantallas y la dualidad mediante
un mapa que transporta campos, graduación, acción y condiciones de
frontera. La relación de supercarga se expresa como igualdad de
operadores con sus dominios, y su teorema de transporte prueba la
conclusión sobre la imagen de la isometría declarada.

El orden expositivo mantiene estas dependencias. Se reproduce primero
el núcleo formal común y su desarrollo, después la generación
coinductiva, el registro \(K\), las dos construcciones de \(\alpha\)
y la incidencia excepcional. Sobre esa base se introducen la
realización hilbertiana, la dualidad, las pantallas y el cociclo de
acción. La repetición del núcleo entre artículos garantiza que este
texto pueda leerse por separado; las fuentes del tratado
integral~\cite{hmtedicionintegral} documentan la procedencia de las
construcciones reunidas. Las demostraciones necesarias se sitúan en
el propio desarrollo y las referencias internas permiten recorrerlas
en ambos sentidos.
